\chapter{Persistence with SQLite}
\chapterintro{The API needs durable state. SQLite gives the project transactions, constraints, indexes, and SQL without introducing a separate database server or an object-relational mapper.}

\begin{objectives}
\begin{itemize}
  \item Translate the Task contract into a constrained relational schema.
  \item Use parameterized statements and explicit transaction boundaries.
  \item Separate database setup and repository operations from HTTP handlers.
  \item Explain per-operation connections, WAL, busy timeout, and migration versioning.
\end{itemize}
\end{objectives}

\section{Why SQLite fits the capstone}

SQLite stores a relational database in a local file and runs in the application
process. Python includes a DB-API interface. The intern can inspect real SQL,
constraints, indexes, transactions, and concurrency behavior without installing
or administering a database server.

This convenience has a boundary. SQLite has one file-oriented write path and is
not a drop-in architecture for every distributed, high-write service. Its
strength here is visibility and low operational overhead.

\section{Schema as durable validation}

\codefile[language=SQL]{../code/final/task_api/schema.sql}

The primary key assigns identity. NOT NULL and CHECK constraints protect required
values even if a future handler forgets validation. Enum checks mirror the API
allow-lists. Length checks provide a final bound. Indexes support filters by
status, priority, and due date.

The schema stores dates and UTC timestamps as consistently formatted text.
Lexicographic ordering works for the chosen ISO forms. SQLite does not enforce a
native date type here, so application validation remains necessary.

\code{PRAGMA user\_version = 1} records schema version. Startup accepts a new
version-zero database and initializes it, accepts version one, and refuses an
unknown version. A real migration system would apply ordered, reversible or
carefully forward-only changes and back up data as policy requires.

\section{Parameterized SQL}

Values never enter SQL through string interpolation. Statements contain
\code{?} placeholders and pass values separately. The driver encodes each value
as data rather than executable SQL syntax.

Dynamic query structure still requires care. The list operation builds WHERE
fragments only from code-owned strings and puts every client value in parameters.
The update operation chooses column names from a fixed allow-list. Parameter
binding cannot substitute table or column identifiers, so those must never come
directly from a request.

\begin{securityboundary}
Escaping quotes manually is not a substitute for parameterization. It is easy to
miss data types, encodings, wildcard semantics, or a future code path. Keep SQL
structure controlled by code and values controlled by parameters.
\end{securityboundary}

\section{Connections and transactions}

\code{Database.connect} creates a connection, configures rows by name, enables
foreign-key enforcement for future schemas, and sets a five-second busy timeout.
Each repository operation opens a connection in a context manager. Successful
write blocks commit; exceptions roll back and close.

Connections are not shared between worker threads. This avoids accidental
cross-thread transaction state. The overhead is acceptable for a local teaching
service and the ownership is obvious.

WAL mode lets readers coexist better with a writer by appending changes to a
write-ahead log before checkpointing. It does not permit unlimited simultaneous
writes. The busy timeout allows short contention to resolve instead of failing
immediately.

\section{The database and repository code}

\codefile{../code/final/task_api/database.py}

The repository returns plain dictionaries matching the API representation. It
does not import HTTP types. A missing row becomes \code{None} or \code{False};
the handler decides that this means a 404 response.

Create uses one INSERT then reads the committed representation. Get and delete
are direct identifier operations. Replace names every writable column. Partial
update builds assignments only from validated writable keys and always adds
\code{updated\_at}.

\section{Filtering, searching, and pagination}

The list method accumulates code-owned predicates and parameters. Search treats
\code{\%} and \code{\_} as literal characters by escaping their LIKE wildcard
meaning. It issues one count query and one ordered page query with the same
conditions.

Ordering uses \code{created\_at DESC, id DESC}. The identifier breaks timestamp
ties, producing deterministic pages for one database state. The count and page
are separate statements, so a concurrent write can occur between them. A stricter
snapshot would wrap both in one read transaction. This implementation accepts
the small inconsistency and documents it.

\section{Repository tests use temporary files}

An in-memory SQLite database is normally private to one connection. Because the
repository intentionally opens short-lived connections, tests use a temporary
database file instead of \code{:memory:}. Each test case initializes a fresh file
and deletes its temporary directory afterward. This mirrors connection behavior
without sharing state between tests.

\begin{pitfall}
Do not put a production-like SQLite file inside a directory that tests delete,
and do not point tests at the developer's real \code{tasks.db}. Test isolation
is a data-safety property.
\end{pitfall}

\summarytitle

SQLite supplies durable constraints and transactions behind a repository that
does not know HTTP. Parameterized SQL separates code-owned structure from client
values. Per-operation connections simplify thread ownership; WAL and busy timeout
improve modest concurrency. Schema versioning makes compatibility explicit.

\reviewtitle
\begin{enumerate}
  \item Which task rules are enforced both in Python and in SQL, and why?
  \item Why can parameters replace values but not arbitrary column names?
  \item Why do tests use a temporary file instead of a simple in-memory database?
  \item What inconsistency can occur between list count and page queries?
\end{enumerate}

\begin{exercise}[Constraint failure]
Bypass HTTP validation and call the repository with an invalid status. Observe
the database constraint failure and transaction rollback. Explain why the HTTP
layer should still validate first and return a controlled 422.
\end{exercise}

\begin{exercise}[Query plan]
Populate tasks with several statuses and due dates. Use SQLite's query-plan
inspection to see which filters can use indexes. Compare an exact status filter
with a leading-wildcard text search and explain the difference.
\end{exercise}

